
\documentclass[11pt,a4paper]{article}
\usepackage[a4paper,margin=2.0cm]{geometry}
\usepackage[T1]{fontenc}
\usepackage[utf8]{inputenc}
\usepackage{microtype,setspace}
\usepackage{amsmath,amssymb}
\usepackage[super,sort&compress,comma]{natbib}
\usepackage{booktabs,longtable,array,tabularx}
\usepackage{graphicx,xcolor,float,enumitem,hyperref,caption,subcaption}
\setstretch{1.06}
\setlength{\parindent}{0pt}
\setlength{\parskip}{5pt}
\setlength{\emergencystretch}{3em}
\hypersetup{colorlinks=true,linkcolor=blue,citecolor=blue,urlcolor=blue}
\newcommand{\figsi}{figures}
\renewcommand{\thefigure}{S\arabic{figure}}
\renewcommand{\thetable}{S\arabic{table}}

\title{\textbf{Electronic Supplementary Information for: The Validation Split Defines the Discovery Claim in AI-Guided MOF Screening}}
\author{Mehrdad Asgari, Shayan Abaei, Angelo Di Bella,\\ Andre Hadji-Thomas, and Hosein Alimardani}
\date{}

\begin{document}
\maketitle
\begingroup
\sloppy
\tableofcontents
\endgroup
\clearpage
\section[Scope and interpretation]{Scope and interpretation of the Electronic Supplementary Information}
This Electronic Supplementary Information (ESI) documents the data joins, split definitions, metric logic,
ESI figures, post-processing tables, and reproducibility outputs supporting the main manuscript.
The central claim is deliberately bounded: the study quantifies split-induced changes in MOF ML benchmark
conclusions. Topology labels are used as a severe grouped-holdout stress test; the analysis does not claim
that topology alone is the mechanistic determinant of adsorption. Likewise, the adsorption targets are
used to evaluate validation-regime sensitivity, not to claim complete process-level optimality for carbon
capture, methane storage, or separations.


\subsection*{ESI roadmap}
 Sections 2--4 define the required files, targets, descriptor families, and grouped-split construction. Sections 5--8 document supplementary target distributions, prediction-scale diagnostics, metric heat maps, variance diagnostics, and exact held-out-group post-processing. Figure and table numbers are generated by LaTeX; file names may retain internal development labels and should not be used as the authoritative numbering.
Section~8 also explains the reader-facing labels used in Figure~5, provides a translation between raw machine-readable group IDs and chemistry-facing audit labels, and gives a manual audit protocol for tracing a high-error grouped-split result from a source ID to representative descriptor values, CIF inspection, prediction audit values, and a conservative chemical interpretation.


\section{Input files and expected project layout}

The Python pipeline uses ARC--MOF-derived adsorption and descriptor data \citep{Burner2023ARCMOF}. The GitHub repository associated with this study is:

\begin{center}
\url{https://github.com/Asgari-research/mof-chemical-holdouts}
\end{center}

The repository includes \texttt{clean\_data.csv}, a processed ARC--MOF-derived benchmark input table prepared for this study. This file was generated from ARC--MOF source resources through data cleaning, identifier normalization, adsorption-target organization, descriptor preparation, and construction of machine-learning-ready inputs. It should therefore be treated as a derived benchmark table rather than as the original ARC--MOF release.

The included \texttt{clean\_data.csv} file supports rerunning the benchmark from the processed working table. Depending on the workflow stage, additional grouped metadata files may be required locally, including \texttt{geo-clusters.csv}, \texttt{mc-clusters.csv}, \texttt{func-clusters.csv}, and \texttt{flig-clusters.csv}. Optional inputs for descriptor-consistency checks or full regeneration from source include \texttt{geometric\_properties.csv} and \texttt{all\_topology\_lists.csv}. If \texttt{clean\_data.csv} is not used directly, the pipeline can build it from ARC--MOF-derived raw files, including \texttt{geometric\_properties.csv}, \texttt{post\_comb\_vsa-CO2.csv}, and \texttt{methane.csv}, together with the grouped metadata files. All identifiers are normalized before joining.

Raw ARC--MOF source files, complete ARC--MOF structural/database files, full CIF archives, the complete historical prediction tree, and the large fitted-model archive are not included in this portable V4 Overleaf package. Users who need the original ARC--MOF source data should obtain them from the original ARC--MOF providers and comply with the associated licence and citation requirements.

The computational handover contains the surviving analysis code, processed benchmark tables, compact source data, revised figure files, fold metrics, recovered compact historical OOS aggregates, and the reviewer-analysis outputs. Targeted historical OOS prediction files were recovered from the original run and used to construct the included compact aggregates. The complete historical prediction tree and large fitted-model archive are not included in this portable handover; this does not establish that they never existed. Exact historical non-Python package versions were not recovered. Current revision software versions are therefore reported separately and are not presented as the original environment.

\begin{table}[H]
\centering
\caption{\textbf{Target overview from the memory-safe sampled prediction table used for point-level diagnostics.}
Rows denote sampled prediction records per target, not the full 263,744-row master adsorption table.
Units are mmol g$^{-1}$. Fold-level metrics and the recovered compact historical OOS aggregates are
audited separately from this sampled diagnostic table.}
\label{tab:target_overview}
\begin{tabular}{lrrrrrr}
\toprule
Target & Rows & Mean & Median & SD & Min & Max \\
\midrule
CH$_4$ uptake at 5.8 bar & 240,000 & 4.241 & 4.133 & 1.440 & 0.000 & 11.302 \\
CH$_4$ uptake at 65 bar & 240,000 & 15.555 & 15.071 & 6.657 & 0.000 & 51.599 \\
CO$_2$ uptake at 0.015 bar & 240,000 & 0.166 & 0.090 & 0.247 & 0.000 & 4.639 \\
CO$_2$ uptake at 0.15 bar & 240,000 & 1.029 & 0.619 & 1.035 & 0.000 & 8.014 \\
\bottomrule
\end{tabular}
\end{table}



\section{Descriptor families}
The executed raw input dimensions are 8, 20, 1, and 21 for compact geometry, enriched interpretable, topology-only, and geometry-plus-topology, respectively. The topology one-hot width is training-fold dependent. Table~\ref{tab:descriptor_definitions} records the exact executed inputs; $\epsilon=10^{-9}$.

\begin{table}[H]
\centering
\scriptsize
\caption{\textbf{Executed descriptor-family definitions.} The historical implementation defines \texttt{cavity\_window\_gap}=Dif; it was not silently redefined in this revision.}
\label{tab:descriptor_definitions}
\begin{tabularx}{\textwidth}{>{\raggedright\arraybackslash}p{0.15\textwidth}c>{\raggedright\arraybackslash}p{0.28\textwidth}>{\raggedright\arraybackslash}X}
\toprule
Family & Raw dimension & Inputs & Executed engineered definitions \\
\midrule
Compact geometry & 8 & Density, ASA, AVA, AVAf, POAVA, $D_i$, $D_f$, Dif & None \\
Enriched interpretable & 20 & Compact inputs plus NASA, POAVAf, unit-cell volume, vASA, and eight engineered columns & $D_i/(D_f+\epsilon)$; Dif (stored as \texttt{cavity\_window\_gap}); ASA$/(\mathrm{AVA}+\epsilon)$; AVAf$/(\mathrm{Density}+\epsilon)$; $\log(1+\max(D_i,0))$; $\log(1+\max(D_f,0))$; POAVA$\times D_f$; Density$\times D_f$ \\
Topology only & 1 & Collapsed topology label & Most-frequent imputation and one-hot encoding within the training fold \\
Geometry + topology & 21 & 20 enriched inputs plus topology label & Same numeric definitions plus fold-local topology encoding \\
\bottomrule
\end{tabularx}
\end{table}

\section{Split families and group-size diagnostics}
The split families are random, geometry-grouped, metal-cluster-grouped, functional-group-grouped, linker-grouped, and topology-grouped. Grouped splits hold all members of a group out together. Figure~\ref{fig:si_group_size} reports group-size distributions; this diagnostic should be inspected because group imbalance can make a split intrinsically harsher.

\begin{figure}[H]
\centering
\includegraphics[width=0.78\textwidth]{\figsi/si_v7_group_size_summary_logscale.pdf}

\caption{\textbf{Group-size diagnostic summary for grouped split families.}
(a) Grouping granularity, shown as the number of distinct groups for each grouped split family on a logarithmic scale. (b) Group-size distribution summary, showing median, mean, and maximum group size for each grouped split family on a logarithmic scale. The diagnostic highlights that split severity is affected not only by chemistry or topology, but also by group-count and group-size imbalance across geometry, metal-cluster, functional-family, linker, and topology labels.}

\label{fig:si_group_size}
\end{figure}



The strong imbalance among grouped split families should be considered when interpreting grouped-split severity. In particular, topology grouping is a severe structural-extrapolation test, but it is also affected by topology-label granularity and unequal group sizes.

\begin{table}[H]
\centering
\scriptsize
\caption{\textbf{Group construction, coverage, and original singleton filtering.} Labelled rows are counted before the submitted benchmark removed groups of size one separately within each grouped cohort. Size summaries apply to eligible groups after singleton removal.}
\label{tab:group_construction}
\resizebox{\textwidth}{!}{%
\begin{tabular}{lrrrrrrl}
\toprule
Family & Labelled & Missing & Raw groups & Singleton groups & Eligible rows & Eligible groups & Eligible size: median [min,max] \\
\midrule
Geometry & 263,744 & 0 & 146,411 & 78,631 & 185,113 & 67,780 & 2 [2,26] \\
Metal & 256,328 & 7,416 & 2,447 & 826 & 255,502 & 1,621 & 4 [2,35,572] \\
Functional & 214,531 & 49,213 & 43,744 & 16,319 & 198,212 & 27,425 & 3 [2,2,954] \\
Linker & 256,328 & 7,416 & 125,220 & 53,796 & 202,532 & 71,424 & 2 [2,356] \\
Topology & 263,744 & 0 & 47 & 0 & 263,744 & 47 & 175 [52,99,167] \\
\bottomrule
\end{tabular}
}
\end{table}

Raw Crystalnet categories with fewer than 50 structures were mapped to \texttt{other}: 415 of 461 raw categories, comprising 3,898 rows, were collapsed, yielding 47 final topology groups. Missing labels were excluded only from the split family that required them. These counts distinguish metadata availability from the final modelled cohort.





\noindent\textbf{Model fitting, preprocessing, and leakage control.}
The benchmark uses the same model families across split families so that the comparison isolates the validation regime rather than changing the learning method. Ridge regression provides a linear baseline. Random forest (RF) is a nonlinear tree-ensemble baseline and is used as the main-text reference model. Histogram gradient boosting (HGB) is a second nonlinear tree-based model, and the multilayer perceptron (MLP) is included as a simple neural-network baseline. The full model suite is retained in the ESI and exported metric tables.

All preprocessing steps are fitted inside the training fold through scikit-learn pipelines. Numerical descriptors are median-imputed using the training fold only. Categorical descriptors are imputed using the most frequent training-fold category and are then one-hot encoded when needed. Scaling is applied inside the pipeline for models that require it, such as ridge regression and MLP. This fold-local preprocessing is important: if imputation values, scaling parameters, or categorical encodings were learned from the full data table before splitting, information from the held-out group could indirectly enter the training procedure. Missing target values cannot be used for either training or testing. Missing group labels are handled differently from missing descriptors: a row without the required group label is excluded from that particular grouped-split evaluation, whereas missing descriptor values are imputed within the training fold.

The hyperparameter grids are fixed before comparing split families and are applied consistently across random and grouped evaluations. This keeps the benchmark focused on the question asked in the main manuscript: how does the validation split change the discovery claim when data, target, descriptors, model class, preprocessing rules, and metrics are otherwise held constant?

\section{Training--held-out generalization gaps}
The 60 targeted RF refits used the exact recovered historical outer memberships and each fold's saved selected RF parameters; hyperparameter search was not repeated. We define $\Delta R^2=R^2_{\rm train}-R^2_{\rm heldout}$ and positive MAE deterioration as $\mathrm{MAE}_{\rm heldout}-\mathrm{MAE}_{\rm train}$. Table~\ref{tab:generalization_gap} reports unweighted means across five folds/repeats.

\begin{table}[H]
\centering
\scriptsize
\caption{\textbf{Training--held-out RF metrics for CO$_2$ uptake at 0.15 bar.} MAE quantities are in mmol g$^{-1}$.}
\label{tab:generalization_gap}
\begin{tabular}{llrrrrrr}
\toprule
Descriptor & Split & $R^2_{train}$ & $R^2_{held}$ & $\Delta R^2$ & MAE$_{train}$ & MAE$_{held}$ & MAE deterioration \\
\midrule
Compact & Random & 0.826 & 0.723 & 0.102 & 0.238 & 0.302 & 0.064 \\
Compact & Geometry & 0.829 & 0.706 & 0.123 & 0.244 & 0.323 & 0.078 \\
Compact & Metal & 0.835 & 0.627 & 0.208 & 0.230 & 0.362 & 0.132 \\
Compact & Functional & 0.842 & 0.718 & 0.124 & 0.224 & 0.299 & 0.075 \\
Compact & Linker & 0.842 & 0.719 & 0.123 & 0.230 & 0.309 & 0.078 \\
Compact & Topology & 0.742 & 0.252 & 0.490 & 0.290 & 0.471 & 0.181 \\
Enriched & Random & 0.856 & 0.758 & 0.098 & 0.214 & 0.280 & 0.066 \\
Enriched & Geometry & 0.858 & 0.739 & 0.119 & 0.221 & 0.301 & 0.081 \\
Enriched & Metal & 0.865 & 0.649 & 0.216 & 0.207 & 0.349 & 0.142 \\
Enriched & Functional & 0.870 & 0.748 & 0.121 & 0.202 & 0.280 & 0.078 \\
Enriched & Linker & 0.871 & 0.751 & 0.120 & 0.206 & 0.288 & 0.081 \\
Enriched & Topology & 0.855 & 0.265 & 0.590 & 0.208 & 0.464 & 0.256 \\
\bottomrule
\end{tabular}
\end{table}

Most held-out metrics were closely reproduced in the current environment, but the exercise is not an exact historical reproduction because the original non-Python package versions were not recovered. The largest absolute held-out $R^2$ difference was 0.011318 for enriched topology fold~2; the corresponding absolute MAE difference was 0.000126 mmol g$^{-1}$, with smaller rank differences. No refit was retuned to improve agreement.

\section{Fixed-universe sensitivity}
The common universe contains exactly 214,531 unique structures with non-null geometry, metal, functional, linker, and topology labels. Every split family began from these same IDs. Singleton groups were deliberately retained in this sensitivity so that split-specific pruning could not change the candidate population; this is intentionally different from the submitted grouped benchmark.

\begin{table}[H]
\centering
\scriptsize
\caption{\textbf{Enriched-RF sensitivity on the fixed 214,531-structure universe.} Values are unweighted mean $\pm$ SD across five outer folds/repeats; recovery is shown as a percentage.}
\label{tab:fixed_universe}
\begin{tabular}{lrrrr}
\toprule
Split & $R^2$ & MAE & Spearman & Top-5\% recovery \\
\midrule
Random & $0.764\pm0.001$ & $0.272\pm0.001$ & $0.900\pm0.000$ & $62.1\pm0.8\%$ \\
Geometry & $0.753\pm0.002$ & $0.277\pm0.002$ & $0.898\pm0.001$ & $61.4\pm1.5\%$ \\
Functional & $0.745\pm0.009$ & $0.281\pm0.005$ & $0.896\pm0.005$ & $60.5\pm1.2\%$ \\
Linker & $0.747\pm0.003$ & $0.280\pm0.004$ & $0.896\pm0.003$ & $60.9\pm0.9\%$ \\
Metal & $0.647\pm0.102$ & $0.339\pm0.091$ & $0.833\pm0.036$ & $53.6\pm9.3\%$ \\
Topology & $0.227\pm0.292$ & $0.463\pm0.299$ & $0.685\pm0.069$ & $37.4\pm18.9\%$ \\
\bottomrule
\end{tabular}
\end{table}

The hierarchy is preserved with modest quantitative changes: random, geometry, functional, and linker validation remain relatively favourable, metal holdout remains harder, and topology holdout remains the strongest failure regime. The large topology fold-level spread is retained explicitly. Candidate-population differences alone therefore do not explain the topology-transfer penalty.

\section[Reproducibility and aggregation]{Reproducibility and exact OOS aggregation}
All historical predictions used for recovery and exact list comparisons were strictly outer-OOS. Benchmark $R^2$, MAE, Spearman, Kendall, recovery, and enrichment were computed separately within each held-out fold/repeat and summarized as the unweighted arithmetic mean of five values. Exact cross-split list analysis instead concatenated outer-OOS rows, averaged repeated predictions by \texttt{filename\_norm}, intersected the two split-family candidate tables, and then calculated ranks, top-$k$, Spearman, and Jaccard on that common universe. Random structures contributed one to five OOS predictions when observed; grouped structures generally contributed one. The five random elite sets were not unioned or voted. The random OOS union contained 177,411 structures. Recovered historical aggregates independently reproduced random--topology Spearman 0.917090, Top-5\% Jaccard 0.199675, and Top-1\% Jaccard 0.066066 to floating-point precision.

The shared random seed was 42. Historical outer random validation used five \texttt{ShuffleSplit} repeats with a 20\% test fraction; grouped validation used five \texttt{GroupKFold} folds. Inner selection used three-fold shuffled \texttt{KFold} for random splits and three-fold group-aware \texttt{GroupKFold} for grouped splits. RF fixed settings were bootstrap sampling, \texttt{max\_samples}=0.70, \texttt{max\_features}=0.70, \texttt{min\_samples\_leaf}=5, and one worker per fit; the grid was \texttt{n\_estimators}$\in\{50,80\}$ and \texttt{max\_depth}$\in\{12,20\}$. RF numerical inputs were median-imputed within each training fold. Rare topology values were collapsed before splitting, and categorical inputs were most-frequent-imputed and one-hot encoded with unseen categories ignored.

The original run recorded Python 3.11.14, but exact historical NumPy, pandas, scikit-learn, and SciPy versions were not recovered. The revision environment was Python 3.13.13, NumPy 1.26.4, pandas 2.3.2, and scikit-learn 1.9.0 on Windows 11. These current versions are reported for the new calculations and are not presented as the historical environment. Exact historical assignments for the targeted generalization refits were taken from recovered OOS fold/repeat labels. The complete historical fitted-estimator archive is not included in this portable handover; the complete original environment lock file was not recovered.

\section[Target and prediction-scale diagnostics]{Target distributions and prediction-scale diagnostics}
Figure~\ref{fig:si_target_dist} shows the target distributions used in the benchmark. Figure~\ref{fig:si_pred_scatter} gives representative prediction scatter panels, and Figure~\ref{fig:si_pld_decile} summarizes PLD-decile error structure. These figures help distinguish target-scale effects from split-family effects.

\begin{figure}[H]
\centering
\includegraphics[width=0.94\textwidth]{\figsi/si_v8_target_distributions_clear_labels.pdf}
\caption{\textbf{Target-distribution diagnostics.} Distributions are shown for CO$_2$ uptake at 0.015 bar, CO$_2$ uptake at 0.15 bar, CH$_4$ uptake at 5.8 bar, and CH$_4$ uptake at 65 bar. Uptake values are reported in mmol g$^{-1}$. These histograms use the memory-safe sampled prediction table for point-level diagnostics and should be interpreted separately from the full 263,744-row master adsorption table.}
\label{fig:si_target_dist}
\end{figure}

\begin{figure}[H]
\centering
\includegraphics[width=0.94\textwidth]{\figsi/si_v8_prediction_scatter_panels_clear_labels.pdf}
\caption{\textbf{Representative prediction scatter panels.} Sampled prediction--target scatter plots provide a visual diagnostic of calibration, bias, and outlier structure for selected target/split/model settings. The panels are representative examples used for point-level inspection and are not intended to replace the full fold-level benchmark metrics or exact prediction-file post-processing.}

\label{fig:si_pred_scatter}
\end{figure}

\begin{figure}[H]
\centering
\includegraphics[width=0.94\textwidth]{\figsi/si_v8_pld_decile_errors_clear_labels.pdf}
\caption{\textbf{Prediction errors by PLD decile.} Mean absolute error is summarized across pore-limiting-diameter deciles for the four adsorption targets. The figure is a diagnostic for whether prediction error varies with pore-size regime; it should be interpreted alongside split-family, descriptor, and target-scale effects rather than as a causal attribution of error to PLD alone.}

\label{fig:si_pld_decile}
\end{figure}

\section{Metric heat maps}
Figures~\ref{fig:si_heat_r2}--\ref{fig:si_heat_enrichment} provide aggregated heat-map diagnostics for $R^2$, MAE, Spearman $\rho$, Kendall $\tau$,
top-5\% recovery, and top-5\% enrichment. The displayed heat maps summarize mean metric values by
split family and model after aggregation over the benchmark settings used for each metric. They are SI
counterparts to the main-text split-family and screening figures and should be interpreted as compact
summary views rather than replacements for the full benchmark tables.


\begin{figure}[H]
\centering
\includegraphics[width=0.92\textwidth]{\figsi/si_heatmap_r2.pdf}
\caption{\textbf{Aggregated heat map of mean $R^2$.} Cell values show mean $R^2$ summarized by split family and model after aggregation over the benchmark settings included in the summary table. The heat map is intended as a compact SI overview; target- and descriptor-specific values are provided in the exported benchmark tables.}

\label{fig:si_heat_r2}
\end{figure}

\begin{figure}[H]
\centering
\includegraphics[width=0.92\textwidth]{\figsi/si_heatmap_mae.pdf}
\caption{\textbf{Aggregated heat map of mean absolute error.} Cell values show mean MAE summarized by split family and model after aggregation over the benchmark settings included in the summary table. MAE values are in mmol g$^{-1}$ where the underlying target scale is an adsorption uptake. Target-specific interpretation should account for the different uptake scales of CO$_2$ and CH$_4$ targets.}


\label{fig:si_heat_mae}
\end{figure}

\begin{figure}[H]
\centering
\includegraphics[width=0.92\textwidth]{\figsi/si_heatmap_spearman_rho.pdf}
\caption{\textbf{Aggregated heat map of Spearman rank correlation.} Cell values show mean Spearman $\rho$ summarized by split family and model after aggregation over the benchmark settings included in the summary table. The figure provides a compact view of global rank-order performance and should be interpreted alongside shortlist-level metrics.}

\label{fig:si_heat_spearman}
\end{figure}

\begin{figure}[H]
\centering
\includegraphics[width=0.92\textwidth]{\figsi/si_heatmap_kendall_tau.pdf}
\caption{\textbf{Aggregated heat map of Kendall rank correlation.} Cell values show mean Kendall $\tau$ summarized by split family and model after aggregation over the benchmark settings included in the summary table. The figure provides a complementary rank-correlation diagnostic to Spearman $\rho$.}


\label{fig:si_heat_kendall}
\end{figure}

\begin{figure}[H]
\centering
\includegraphics[width=0.92\textwidth]{\figsi/si_heatmap_top_5pct_overlap.pdf}
\caption{\textbf{Aggregated heat map of top-5\% recovery.} Cell values show mean top-5\% recovery summarized by split family and model after aggregation over the benchmark settings included in the summary table. Here top-5\% recovery refers to recovery of true high-uptake candidates by the predicted top-5\% set, not the predicted-list Jaccard overlap used in the exact elite-list analysis.}

\label{fig:si_heat_overlap}
\end{figure}

\begin{figure}[H]
\centering
\includegraphics[width=0.92\textwidth]{\figsi/si_heatmap_top_5pct_enrichment.pdf}
\caption{\textbf{Aggregated heat map of top-5\% enrichment.} Cell values show mean enrichment of the predicted top-5\% set relative to random selection, summarized by split family and model after aggregation over the benchmark settings included in the summary table. Values greater than one indicate better-than-random recovery of high-uptake candidates.}

\label{fig:si_heat_enrichment}
\end{figure}

\section{Screening-rank and variance diagnostics}
Figure~\ref{fig:si_screen_rank} relates prediction error to screening rank. Figure~\ref{fig:si_variance} reports marginal variance decomposition across split family, descriptor family, model, and target.

\begin{figure}[H]
\centering
\includegraphics[width=0.88\textwidth]{\figsi/si_v7_error_vs_screening_rank_clean.pdf}
\caption{\textbf{Error versus screening-quality diagnostic.} Each point summarizes a benchmark configuration by mean MAE and top-5\% recovery. The figure tests whether configurations with low average error also preserve high-uptake candidate recovery. Points in the lower-right region indicate configurations with both high error and weak shortlist recovery.}

\label{fig:si_screen_rank}
\end{figure}

\begin{figure}[H]
\centering
\includegraphics[width=0.82\textwidth]{\figsi/si_variance_decomposition.pdf}
\caption{\textbf{Variance decomposition across response metrics.} Eta-squared values summarize the marginal contribution of benchmark factors to each response metric. The decomposition is metric-dependent: split family and descriptor family contribute comparably for $R^2$, target identity dominates MAE and RMSE, and descriptor family contributes most strongly to the top-5\% screening metrics in this summary. The result supports split-family reporting without implying that split family is the dominant factor for every metric.}


\label{fig:si_variance}
\end{figure}
Table~S2 shows that benchmark sensitivity is metric-dependent. For $R^2$, split family and descriptor
family contribute comparable marginal variance, supporting the main-text claim that validation regime
can materially change accuracy-based conclusions. For MAE and RMSE, target identity contributes the
largest variance component, reflecting the different physical scales of the four adsorption targets. For
top-5\% recovery and enrichment, descriptor family contributes the largest marginal component in this
summary. Thus, the variance decomposition supports the importance of split-family reporting, but it
does not imply that split family is the dominant factor for every metric.


\begin{table}[H]
\centering
\caption{Variance-decomposition summary. Eta-squared values are computed from pipeline summary tables.}
\label{tab:si_variance}
\small
\begin{tabular}{llr}
\toprule
Response & Factor & $\eta^2$ \\
\midrule
r2 & split family & 0.183 \\
r2 & descriptor family & 0.186 \\
r2 & model & 0.011 \\
r2 & target slug & 0.099 \\
mae & split family & 0.011 \\
mae & descriptor family & 0.177 \\
mae & model & 0.003 \\
mae & target slug & 0.361 \\
rmse & split family & 0.008 \\
rmse & descriptor family & 0.180 \\
rmse & model & 0.003 \\
rmse & target slug & 0.364 \\
top 5pct overlap & split family & 0.036 \\
top 5pct overlap & descriptor family & 0.465 \\
top 5pct overlap & model & 0.044 \\
top 5pct overlap & target slug & 0.263 \\
top 5pct enrichment & split family & 0.036 \\
top 5pct enrichment & descriptor family & 0.465 \\
top 5pct enrichment & model & 0.044 \\
top 5pct enrichment & target slug & 0.263 \\
\bottomrule
\end{tabular}
\end{table}


\section{Held-out groups and exact post-processing}
The pipeline writes full held-out prediction files as per-experiment CSV files and computes group-resolved
errors by streaming those files from disk. Each row in the high-error-group analysis represents a held-out
group summarized over multiple structures, not an individual MOF. The plotted value in Figure~5 of the
main text is the group-level mean absolute error. Table~\ref{tab:si_label_translation} explains how
machine-readable IDs such as \texttt{metal:2157.0}, \texttt{func:18.0}, and \texttt{geo:94288} are translated
into reader-facing audit labels without changing the underlying numerical data.





The representative structural-audit rows discussed in the main manuscript are provided in the deposited source-data package as \texttt{table\_case\_study\_candidates.csv}. This file contains the machine-readable identifiers, split/model/descriptor setting, true and predicted uptake values, error summaries, rank-shift information where available, and descriptor values used to construct the diagnostic structural-audit examples. The file is provided for traceability only; the listed rows are representative audit cases and should not be interpreted as best-MOF candidates, process-optimal sorbents, or an independent validation set.



\subsection[Figure 5 labels and source-ID traceability]{Reader-facing labels for Figure 5 and source-ID traceability}
Raw group IDs are essential for reproducibility because they point exactly to the grouped-split metadata
and post-processing tables. They are, however, difficult to read in the main text. The manuscript therefore
uses a two-level convention: the main text and Figure~5 use chemistry-facing audit labels, while the exact
machine-readable source IDs are retained in source-data tables and in Table~\ref{tab:si_label_translation}.
Chemically specific labels are used only when they are supported by filename-encoded tokens in the
cluster files. For example, metal-cluster ID~2157 consistently carries the filename-encoded
Cu$_2$O$_8$N$_2$ and BDC\_A tokens in \texttt{mc-clusters.csv}, so it is described as a
Cu$_2$O$_8$N$_2$--BDC node-family audit group. By contrast, functional-family clusters such as
\texttt{func:18.0} and \texttt{func:6.0} do not contain a plain-language functional motif name in the
uploaded cluster table, so they are conservatively labelled as functional-family labels rather than as
specific named functional motifs. These labels make Figure~5 easier to read while avoiding unsupported
mechanistic or structural claims.


\begin{table}[H]
\centering
\scriptsize
\caption{\textbf{Translation between machine-readable Figure~5 group IDs and reader-facing audit labels.}
The source ID preserves exact traceability to the grouped-split metadata and post-processing tables.
The reader-facing label is used for clarity in the main-text figure. Chemical specificity is introduced
only when supported by filename tokens or cluster-file evidence; otherwise a conservative family-level
label is used.}
\label{tab:si_label_translation}
\setlength{\tabcolsep}{4pt}
\begin{tabular}{p{0.15\textwidth}p{0.25\textwidth}p{0.33\textwidth}p{0.17\textwidth}}
\toprule
\textbf{Source ID} & \textbf{Reader-facing label} & \textbf{Evidence used for label} & \textbf{Caution} \\
\midrule
\texttt{metal:2157.0} & Cu$_2$O$_8$N$_2$--BDC node-family group & Filename tokens in \texttt{mc-clusters.csv} consistently encode Cu$_2$O$_8$N$_2$ and BDC\_A & Not proof of identical adsorption mechanism \\
\texttt{metal:1858.0} & Cu$_2$O$_8$N$_2$--AZO-rich node-family group & Cu$_2$O$_8$N$_2$ token; AZO\_A-rich cluster membership & Rich label, not necessarily every row identical \\
\texttt{metal:1919.0} & Zn$_2$O$_8$N$_2$--AZO node-family group & Filename tokens encode Zn$_2$O$_8$N$_2$ and AZO\_A & Node-family audit label only \\
\texttt{metal:1837.0} & Zn$_2$O$_8$ node-family group & Filename tokens encode Zn$_2$O$_8$ with mixed organic/framework tokens & Mixed linker/framework contexts \\
\texttt{metal:1120.0} & CSD node-family group 1120 & CSD/DB12-style filenames without a simple node token & Inspect representative CIFs \\
\texttt{metal:2170.0} & Cu$_2$O$_8$N$_2$--IRMOF-rich node-family group & Cu$_2$O$_8$N$_2$ token; IRMOF-rich filename tokens & Family-level label \\
\texttt{metal:1901.0} & Zn$_2$O$_8$N$_2$--IRMOF-rich node-family group & Zn$_2$O$_8$N$_2$ token; IRMOF-rich filename tokens & Family-level label \\
\texttt{metal:313.0} & Zn$_4$O$_{13}$ node-family group & Filename tokens encode Zn$_4$O$_{13}$ & Mixed linker/framework contexts \\
\texttt{metal:1898.0} & Cu$_2$O$_8$N$_2$--DPAC-rich node-family group & Cu$_2$O$_8$N$_2$ token; DPAC\_A-rich filename tokens & Rich label, not necessarily every row identical \\
\texttt{linker:2934.0} & Linker-family group 2934 & Linker-cluster ID from \texttt{flig-clusters.csv} & Exact linker identity should be checked from source structures \\
\texttt{func:18.0} & Functional-family 18 & Functional-cluster ID from \texttt{func-clusters.csv}; no plain-language motif name encoded & Do not invent a named motif \\
\texttt{func:6.0} & Functional-family 6 & Functional-cluster ID from \texttt{func-clusters.csv}; no plain-language motif name encoded & Do not invent a named motif \\
\texttt{geo:*} & Pore-regime group + source ID & Geometry-cluster ID from the grouped metadata & Pore-regime label, not mechanism \\
\bottomrule
\end{tabular}
\end{table}


Table~\ref{tab:si_hard_groups} lists representative high-error held-out groups for
the anchor CO$_2$ uptake target. The complete exact group-resolved error table can be regenerated by
rerunning the script. The table is intended as a diagnostic audit of where prediction error concentrates,
not as a list of optimal sorbents or process-ready candidates.

\begin{table}[H]
\centering
\caption{\textbf{Representative high-error held-out groups for CO$_2$ uptake at 0.15 bar with enriched interpretable descriptors and RF.}
MAE, mean target, and mean prediction are reported in mmol g$^{-1}$. Rows are selected diagnostic
examples from exact post-processing of full prediction files; each row represents a held-out group, not an individual MOF. Reader-facing audit labels are paired with exact source IDs for traceability.}
\label{tab:si_hard_groups}
\scriptsize
\setlength{\tabcolsep}{3pt}
\begin{tabular}{p{0.17\textwidth}p{0.24\textwidth}p{0.16\textwidth}rrrr}
\toprule
Split family & Reader-facing audit label & Source ID & $n$ & MAE & Mean target & Mean prediction \\
\midrule
Topology-grouped & Cu$_2$O$_8$N$_2$--BDC node-family group & \texttt{metal:2157.0} & 25 & 6.911 & 8.182 & 1.271 \\
Metal-cluster-grouped & Cu$_2$O$_8$N$_2$--BDC node-family group & \texttt{metal:2157.0} & 25 & 6.496 & 8.182 & 1.686 \\
Random & Cu$_2$O$_8$N$_2$--BDC node-family group & \texttt{metal:2157.0} & 27 & 5.099 & 8.183 & 3.084 \\
Topology-grouped & CSD node-family group 1120 & \texttt{metal:1120.0} & 31 & 3.719 & 5.429 & 1.710 \\
Functional-grouped & Cu$_2$O$_8$N$_2$--IRMOF-rich node-family group & \texttt{metal:2170.0} & 22 & 3.639 & 5.177 & 1.537 \\
Metal-cluster-grouped & Cu$_2$O$_8$N$_2$--AZO-rich node-family group & \texttt{metal:1858.0} & 22 & 3.615 & 4.192 & 0.576 \\
Metal-cluster-grouped & Pore-regime group 94288 & \texttt{geo:94288} & 26 & 3.319 & 4.748 & 1.429 \\
Topology-grouped & Pore-regime group 94288 & \texttt{geo:94288} & 26 & 3.229 & 4.748 & 1.519 \\
Random & AZO-rich functional-family group 6004 & \texttt{func:6004.0} & 23 & 3.116 & 4.196 & 1.080 \\
Topology-grouped & Pore-regime group 70199 & \texttt{geo:70199} & 26 & 3.090 & 5.137 & 2.123 \\
Metal-cluster-grouped & Pore-regime group 109675 & \texttt{geo:109675} & 20 & 3.071 & 4.243 & 1.172 \\
Metal-cluster-grouped & Linker-family group 14604 & \texttt{linker:14604.0} & 20 & 3.071 & 4.243 & 1.172 \\
\bottomrule
\end{tabular}
\end{table}
\noindent\textit{Note.} Some representative high-error rows can share the same support size and summary statistics across different metadata or audit-label families because the same set of structures may be selected under more than one grouping lens. Such rows are retained when they correspond to different split-family or metadata interpretations; the exact source IDs preserve the distinction.

\newpage

The main text includes a structural audit figure based on three representative case-study candidates.
Tables~\ref{tab:structural_audit_descriptors} and \ref{tab:structural_audit_predictions} provide the
numerical source-data summary for these cases. The examples are diagnostic candidates selected from
the grouped-extrapolation audit and should not be interpreted as optimal sorbents or process-screening
recommendations.

\begin{table}[H]
\centering
\caption{\textbf{Descriptor and group-label summary for the main-text structural audit cases.}
Descriptor values correspond to the three cases shown in the main-text structural audit figure. Case A belongs to the Cu$_2$O$_8$N$_2$--BDC node-family audit group, retained as metal-cluster ID~2157. These
cases are diagnostic examples from grouped-extrapolation audits, not best-MOF or process-screening
recommendations.}
\label{tab:structural_audit_descriptors}
\small
\setlength{\tabcolsep}{4pt}
\resizebox{\textwidth}{!}{%
\begin{tabular}{@{}p{0.05\textwidth}p{0.22\textwidth}p{0.16\textwidth}p{0.09\textwidth}p{0.10\textwidth}rrrr@{}}
\toprule
Case & Framework label & Split & Topology & Metal cluster & Df/PLD & Di/LCD & Density & AVAf \\
 &  &  &  &  & (\AA) & (\AA) & (g cm$^{-3}$) &  \\
\midrule
A & pcu Cu framework; Cu$_2$O$_8$N$_2$--BDC node-family group & Topology-grouped & pcu & 2157.0 & 7.252 & 7.748 & 0.965 & 0.112 \\
B & dia Cu framework & Geometry-grouped & dia & 108.0 & 8.351 & 8.712 & 1.035 & 0.171 \\
C & large-pore tbo framework & Topology-grouped & tbo & 1866.0 & 12.399 & 24.182 & 0.309 & 0.524 \\
\bottomrule
\end{tabular}
}
\end{table}

\begin{table}[H]
\centering
\caption{\textbf{Prediction-audit summary for the main-text structural audit cases.}
Uptake and error values are reported in mmol g$^{-1}$. ASA is reported in m$^2$ g$^{-1}$. POAVA is
reported using the same units as the source descriptor table. The examples are intended as
structure-level diagnostics of split-induced prediction error, not as material recommendations.}
\label{tab:structural_audit_predictions}
\small
\setlength{\tabcolsep}{4pt}
\begin{tabular}{@{}p{0.05\textwidth}rrrrrrrr@{}}
\toprule
Case & ASA & POAVA & True & Random pred. & Grouped pred. & Random error & Grouped error & $\Delta$ error \\
 & (m$^2$ g$^{-1}$) &  &  &  &  &  &  &  \\
\midrule
A & 306.3  & 851.8   & 8.09 & 4.85 & 1.58 & 3.239 & 6.509 & 3.270 \\
B & 322.5  & 976.7   & 6.18 & 4.92 & 1.74 & 1.255 & 4.443 & 3.188 \\
C & 4690.9 & 21111.7 & 3.77 & 3.37 & 0.39 & 0.401 & 3.386 & 2.985 \\
\bottomrule
\end{tabular}
\end{table}

\subsection[Structural-audit interpretation boundaries]{Interpretation boundaries for structural-audit cases}
\label{sec:si_structural_audit_boundaries}

The structural-audit cases are a post hoc interpretability layer applied to benchmark outputs that have already been generated by the split-strategy pipeline. They are not an additional validation set and do not introduce new model training, new descriptor calculation, or new database joins. Their role is to connect high-error grouped-split regimes to chemically inspectable structures while preserving traceability to the exported source data.

Each audit case links four levels of evidence: the machine-readable source identifier and reader-facing audit label; the support size and error statistics for the relevant held-out group; the descriptor and topology context of representative structures; and the prediction-level comparison between random-split and grouped-split outputs. This structure allows the reader to trace how an aggregate high-error regime becomes a chemically interpretable example without treating the example as a material recommendation.

The resulting interpretations should therefore be read conservatively. A structural-audit case can illustrate node-family extrapolation, pore-regime transfer, topology/connectivity transfer, linker-family transfer, functional-motif transfer, high-uptake underprediction, or shortlist dropout. It does not by itself establish a microscopic adsorption mechanism, prove that topology alone controls uptake, or identify a process-optimal sorbent. Such mechanistic or process-level conclusions would require additional molecular simulation, site-level analysis, experimental validation, or process modelling.




\begin{table}[H]
\centering
\small
\setlength{\tabcolsep}{3pt}
\caption{\textbf{Source-data fields used for structural-audit traceability.}
This table defines the fields used to convert machine-readable grouped-split outputs into conservative chemical annotations. It is a protocol table rather than a new benchmark result. Group labels should be retained exactly for traceability, and chemical descriptions should be filled only after inspecting representative source-data rows and CIF-derived structures.}
\begin{tabular}{p{0.16\textwidth} p{0.18\textwidth} p{0.18\textwidth} p{0.20\textwidth} p{0.20\textwidth}}
\toprule
\textbf{Field} & \textbf{What to record} & \textbf{Example/source} & \textbf{Use in interpretation} & \textbf{Caution} \\
\midrule
Group label & Exact machine-readable label plus reader-facing audit label & \texttt{metal:2157.0}; Cu$_2$O$_8$N$_2$--BDC node-family group & Preserves traceability while remaining readable & Do not infer chemistry from ID alone \\
Support/error & Number of structures, MAE, target mean, prediction mean & ESI hard-group table or source-data CSV & Identifies high-error regimes with adequate support & Avoid overinterpreting very small groups \\
Structure context & Topology, node/metal context, linker appearance, pore channel & CIF rendering or representative CIF & Provides a chemically inspectable case & Visual inspection is not a mechanism proof \\
Descriptor context & PLD/Df, LCD/Di, density, AVAf, ASA, pore volume & Structural audit source-data table & Connects error to pore-regime and geometry descriptors & Descriptor values are summaries, not full adsorption physics \\
Prediction audit & True uptake, random prediction, grouped prediction, errors, rank shift & Prediction-level audit table & Distinguishes underprediction, overprediction, or shortlist dropout & Do not treat a single case as a new validation set \\
Audit category & Node-family, pore-regime, topology/connectivity, linker-family, motif, or shortlist-dropout case & Manual annotation after inspection & Gives readers a concise interpretation route & Use conservative language: ``illustrates'' or ``is consistent with'' \\
\bottomrule
\end{tabular}
\label{tab:manual_audit_worksheet}
\end{table}

\clearpage
\bibliographystyle{rsc}
\bibliography{references}


\end{document}
